% lhopital-limits.tex — L'Hôpital limits for TKF parameters as λ → μ
% To be \input from tkf.tex or read standalone

\section{Equal-Rate Limits for TKF Parameters}
\label{sec:lhopital}

For biologically realistic equilibrium lengths under TKF91/TKF92, one is typically driven into the near-critical regime
$\insrate \approx \delrate$ ($\kappa \approx 1$).
Under TKF91, the expected equilibrium sequence length is
$\kappa/(1-\kappa)$ characters.
Under TKF92 with fragment extension probability $\ext$,
it is $\kappa/((1-\ext)(1-\kappa))$ characters (over $\kappa/(1-\kappa)$ fragments).
The average protein in UniProtKB/Swiss-Prot is approximately
360 amino acids long \cite{UniProtStatistics}.
Even under TKF92 with a mean indel length of $1/(1-\ext) \approx 3$
residues (matching the empirically observed predominance of
1--5 residue indels in protein evolution \cite{Cartwright2009,Ajawatanawong2013}),
the required $\kappa$ is $120/121 \approx 0.992$.
Under TKF91, $\kappa = 360/361 \approx 0.997$.
For DNA sequences (a typical prokaryotic ORF or eukaryotic
intron-free coding sequence is $\sim$1000 bp
\cite{XuEtAl2006}; intron-bearing eukaryotic genes and syntenic
genomic regions can be much longer), $\kappa > 0.999$.
Since $1 - \kappa = (\delrate - \insrate)/\delrate$, these
biological sequence lengths force $\insrate$ very close to $\delrate$.

The TKF91 transition probabilities $\alpha, \beta, \gamma$ and the
BDI score/sufficient statistics all contain factors of
$(\delrate - \insrate)$ in denominators.
When $\insrate \to \delrate$ ($\kappa \to 1$),
these become $0/0$ indeterminate forms.
This section derives the L'H\^opital limits, providing
numerically stable expressions in the equal-rate regime.

Let $\varepsilon = \delrate - \insrate \to 0$ and $s = \delrate\,\evoltime$.
We parameterize $\insrate = \delrate - \varepsilon$ and take
$\varepsilon \to 0^+$.

Throughout, we write $\Phi = 1 - \alpha$ for the complement of the
survival probability ($\Phi$ is well-defined at any $(\insrate,\delrate)$).
In the equal-rate regime, where $s = \delrate\evoltime$ and $\alpha = e^{-s}$,
we additionally define:
\[
\phi = \Phi/s = (1-e^{-s})/s
\]
Note $\phi(0) = 1$ and $\phi(s) \to 0$ as $s \to \infty$.

\subsection{TKF91 Transition Parameters}

\paragraph{$\alpha$ (survival probability).}
$\alpha = e^{-\delrate\evoltime}$ has no dependence on $\insrate$;
no singularity.

\paragraph{$\beta$ (insertion continuation probability).}
\begin{equation}
\beta = \frac{\insrate(e^{-\insrate\evoltime} - e^{-\delrate\evoltime})}
             {\delrate\,e^{-\insrate\evoltime} - \insrate\,e^{-\delrate\evoltime}}
\end{equation}

Write $e^{-\insrate\evoltime} = \alpha\,e^{\varepsilon\evoltime}$ and
let $u = \varepsilon\evoltime$.
Then:
\begin{align*}
\text{Numerator}
  &= (\delrate - \varepsilon)\,\alpha\,(e^u - 1) \\
\text{Denominator}
  &= \alpha\bigl(\delrate\,e^u - \delrate + \varepsilon\bigr)
\end{align*}
Cancelling $\alpha$ and expanding $e^u = 1 + u + u^2/2 + \cdots$
with $u = \varepsilon\evoltime$:
\begin{align*}
\text{Numer} &= (\delrate - \varepsilon)(\varepsilon\evoltime
  + \varepsilon^2\evoltime^2/2 + \cdots)
  = \varepsilon\bigl(\delrate\evoltime + O(\varepsilon)\bigr) \\
\text{Denom} &= \delrate\varepsilon\evoltime
  + \varepsilon(1 + O(\varepsilon))
  = \varepsilon\bigl(\delrate\evoltime + 1 + O(\varepsilon)\bigr)
\end{align*}
After cancelling $\varepsilon$:
\begin{equation}
\boxed{\beta \;\xrightarrow{\;\insrate\to\delrate\;}\;
  \frac{s}{1+s}}
\label{eq:beta-limit}
\end{equation}

\paragraph{$\gamma$ (orphan insertion probability).}
Using $\gamma = 1 - \delrate\beta/(\insrate(1-\alpha))$:
\begin{equation}
\boxed{\gamma \;\xrightarrow{\;\insrate\to\delrate\;}\;
  1 - \frac{1}{(1+s)\,\phi}
  = \frac{(1+s)\phi - 1}{(1+s)\phi}}
\label{eq:gamma-limit}
\end{equation}

\paragraph{Complementary parameters.}
\begin{align}
1-\beta &\;\to\; \frac{1}{1+s} \label{eq:1mbeta-limit}\\
1-\gamma &\;\to\; \frac{1}{(1+s)\,\phi} \label{eq:1mgamma-limit}
\end{align}


%%%%%%%%%%%%%%%%%%%%%%%%%%%%%%%%%%%%%%%%%%%%%%%%%%%%%%%%%%%%%%%%%%%%%%%%
\subsection{Score Derivatives: General Case}
\label{sec:score-general}

Define $\eta = e^{-\insrate\evoltime}$ and
$\delta = \delrate\eta - \insrate\alpha$
(the denominator of $\beta$).
The score of the observed log-likelihood
decomposes into 8 log-parameter groups.
Below we list the derivatives of each log-factor with respect to
$\insrate$ and $\delrate$, from which the log-elasticities follow
by multiplication:
\[
\frac{\partial\log\xi}{\partial\log\insrate}
= \insrate\,\frac{\partial\log\xi}{\partial\insrate}, \qquad
\frac{\partial\log\xi}{\partial\log\delrate}
= \delrate\,\frac{\partial\log\xi}{\partial\delrate}, \qquad
\frac{\partial\log\xi}{\partial\log(\insrate{+}\delrate)}
= \frac{\partial\log\xi}{\partial\log\insrate}
+ \frac{\partial\log\xi}{\partial\log\delrate}
\]
The third identity holds because differentiating with respect to
$\log(\insrate{+}\delrate)$ at fixed ratio $\kappa = \insrate/\delrate$
is equivalent to scaling both rates by a common factor $c$ and
differentiating with respect to $\log c$, which gives
$\insrate\,\partial_\insrate + \delrate\,\partial_\delrate$.

\subsubsection{\texorpdfstring{$\log\alpha$ and $\log(1{-}\alpha)$}{log alpha and log(1-alpha)}}

Since $\alpha = e^{-\delrate\evoltime}$ depends only on $\delrate$:
\begin{align}
\frac{\partial\log\alpha}{\partial\insrate} &= 0, &
\frac{\partial\log\alpha}{\partial\delrate} &= -\evoltime
\label{eq:dlogalpha}
\\[4pt]
\frac{\partial\log(1{-}\alpha)}{\partial\insrate} &= 0, &
\frac{\partial\log(1{-}\alpha)}{\partial\delrate}
  &= \frac{\evoltime\alpha}{1-\alpha}
  = \frac{\evoltime\alpha}{\Phi}
\label{eq:dlog1malpha}
\end{align}

\subsubsection{\texorpdfstring{$\log\beta$ and $\log(1{-}\beta)$}{log beta and log(1-beta)}}

From $\log\beta = \log\insrate + \log(\eta - \alpha) - \log\delta$:
\begin{align}
\frac{\partial\log\beta}{\partial\insrate}
  &= \frac{1}{\insrate}
     - \frac{\evoltime\eta}{\eta - \alpha}
     + \frac{\evoltime\delrate\eta + \alpha}{\delta}
\label{eq:dlogbeta-dl}
\\[4pt]
\frac{\partial\log\beta}{\partial\delrate}
  &= \frac{\evoltime\alpha}{\eta - \alpha}
     - \frac{\eta + \evoltime\insrate\alpha}{\delta}
\label{eq:dlogbeta-dm}
\end{align}

Both $(\eta - \alpha) \to 0$ and $\delta \to 0$ as $\insrate \to \delrate$,
so each term individually diverges; the combination is finite
(Section~\ref{sec:score-limits}).

For $\log(1{-}\beta)$, using the chain rule
$\partial\log(1{-}\beta) = -\frac{\beta}{1-\beta}\,\partial\log\beta$:
\begin{align}
\frac{\partial\log(1{-}\beta)}{\partial\insrate}
  &= -\frac{\beta}{1-\beta}
     \cdot\frac{\partial\log\beta}{\partial\insrate}
\label{eq:dlog1mbeta-dl}
\\[4pt]
\frac{\partial\log(1{-}\beta)}{\partial\delrate}
  &= -\frac{\beta}{1-\beta}
     \cdot\frac{\partial\log\beta}{\partial\delrate}
\label{eq:dlog1mbeta-dm}
\end{align}

\subsubsection{\texorpdfstring{$\log(1{-}\gamma)$ and $\log\gamma$}{log(1-gamma) and log gamma}}

From $1 - \gamma = \delrate\beta/(\insrate(1{-}\alpha))$,
so $\log(1{-}\gamma) = \log\delrate + \log\beta - \log\insrate - \log(1{-}\alpha)$:
\begin{align}
\frac{\partial\log(1{-}\gamma)}{\partial\insrate}
  &= \frac{\partial\log\beta}{\partial\insrate} - \frac{1}{\insrate}
\label{eq:dlog1mgamma-dl}
\\[4pt]
\frac{\partial\log(1{-}\gamma)}{\partial\delrate}
  &= \frac{1}{\delrate}
     + \frac{\partial\log\beta}{\partial\delrate}
     - \frac{\evoltime\alpha}{1-\alpha}
\label{eq:dlog1mgamma-dm}
\end{align}

For $\log\gamma$:
\begin{align}
\frac{\partial\log\gamma}{\partial\insrate}
  &= -\frac{1-\gamma}{\gamma}
     \cdot\frac{\partial\log(1{-}\gamma)}{\partial\insrate}
\label{eq:dloggamma-dl}
\\[4pt]
\frac{\partial\log\gamma}{\partial\delrate}
  &= -\frac{1-\gamma}{\gamma}
     \cdot\frac{\partial\log(1{-}\gamma)}{\partial\delrate}
\label{eq:dloggamma-dm}
\end{align}


\subsubsection{Summary Table: General Case}
\label{sec:score-table-general}

\begin{center}
\renewcommand{\arraystretch}{1.6}
\begin{tabular}{c|cccccc}
& $\log\alpha$ & $\log(1{-}\alpha)$
& $\log\beta$ & $\log(1{-}\beta)$
& $\log\gamma$ & $\log(1{-}\gamma)$ \\
\hline
$\dfrac{\partial}{\partial\log\insrate}$
  & $0$
  & $0$
  & $1 - \dfrac{s_\insrate \eta}{\eta-\alpha}
        + \dfrac{s_\insrate\delrate\eta/\insrate + \insrate\alpha/\insrate}{\delta/\insrate}$
  & $-\dfrac{\beta}{1-\beta}\cdot[\cdot]_\beta^\insrate$
  & $-\dfrac{1-\gamma}{\gamma}\cdot[\cdot]_{1-\gamma}^\insrate$
  & $[\cdot]_\beta^\insrate - 1$
\\[6pt]
$\dfrac{\partial}{\partial\log\delrate}$
  & $-s$
  & $\dfrac{s\alpha}{\Phi}$
  & $\dfrac{s\alpha}{\eta-\alpha} - \dfrac{\delrate\eta + s\insrate\alpha}{\delta}$
  & $-\dfrac{\beta}{1-\beta}\cdot[\cdot]_\beta^\delrate$
  & $-\dfrac{1-\gamma}{\gamma}\cdot[\cdot]_{1-\gamma}^\delrate$
  & $1 + [\cdot]_\beta^\delrate - \dfrac{s\alpha}{\Phi}$
\\[6pt]
$\dfrac{\partial}{\partial\log(\insrate{+}\delrate)}$
  & \multicolumn{6}{c}{(sum of above two rows)}
\end{tabular}
\end{center}

\noindent
where $s_\insrate = \insrate\evoltime$, $s = \delrate\evoltime$,
and $[\cdot]_\xi^\theta$ denotes the $\partial\log\xi/\partial\log\theta$ entry.

The table is dense because each $\log\beta$ and $\log\gamma$ entry
involves the singular terms $(\eta-\alpha)^{-1}$ and $\delta^{-1}$.
These singularities cancel in each entry, as we show next.


%%%%%%%%%%%%%%%%%%%%%%%%%%%%%%%%%%%%%%%%%%%%%%%%%%%%%%%%%%%%%%%%%%%%%%%%
\subsection{Score Derivatives: L'H\^opital Limits}
\label{sec:score-limits}

We systematically expand each derivative to first order in
$\varepsilon = \delrate - \insrate$ and extract the finite limit.
With $u = \varepsilon\evoltime$, the key expansions are:
\begin{align}
\eta - \alpha &= \alpha(e^u - 1)
  = \alpha\bigl(u + u^2/2 + \cdots\bigr)
\label{eq:eta-alpha-expand}
\\
\delta &= \delrate\eta - \insrate\alpha
  = \alpha\varepsilon(1 + s + O(\varepsilon))
\label{eq:delta-expand}
\end{align}
and the Laurent expansion $1/(e^u - 1) = 1/u - 1/2 + u/12 - \cdots$,
giving:
\begin{align}
\frac{\evoltime\eta}{\eta - \alpha}
  = \frac{\evoltime e^u}{e^u - 1}
  &= \frac{1}{\varepsilon} + \frac{\evoltime}{2} + O(\varepsilon)
\label{eq:singular-A}
\\
\frac{\evoltime\alpha}{\eta - \alpha}
  = \frac{\evoltime}{e^u - 1}
  &= \frac{1}{\varepsilon} - \frac{\evoltime}{2} + O(\varepsilon)
\label{eq:singular-B}
\\
\frac{\evoltime\delrate\eta + \alpha}{\delta}
  &= \frac{1}{\varepsilon} + \frac{s\evoltime}{2(1+s)} + O(\varepsilon)
\label{eq:singular-C}
\\
\frac{\eta + \evoltime\insrate\alpha}{\delta}
  &= \frac{1}{\varepsilon} - \frac{s\evoltime}{2(1+s)} + O(\varepsilon)
\label{eq:singular-D}
\end{align}

Derivation of~\eqref{eq:singular-C}: the numerator is
$\evoltime\delrate\eta + \alpha
= \alpha(se^u + 1)
= \alpha\bigl((1{+}s) + s\varepsilon\evoltime + O(\varepsilon^2)\bigr)$
and the denominator is
$\delta = \alpha\varepsilon\bigl((1{+}s) + s\varepsilon\evoltime/2 + O(\varepsilon^2)\bigr)$
from~\eqref{eq:delta-expand}.  So:
\[
\frac{\evoltime\delrate\eta + \alpha}{\delta}
= \frac{(1{+}s) + s\varepsilon\evoltime + \cdots}
       {\varepsilon\bigl((1{+}s) + s\varepsilon\evoltime/2 + \cdots\bigr)}
= \frac{1}{\varepsilon}
  \cdot\frac{1 + \frac{s\varepsilon\evoltime}{1+s} + \cdots}
            {1 + \frac{s\varepsilon\evoltime}{2(1+s)} + \cdots}
\]
The $O(\varepsilon)$ coefficients differ: $s\evoltime/(1{+}s)$ in the
numerator vs.\ $s\evoltime/(2(1{+}s))$ in the denominator.
Their difference gives the $O(1)$ correction:
$1/\varepsilon + s\evoltime/(2(1{+}s)) + O(\varepsilon)$.
Similarly for~\eqref{eq:singular-D}.

\subsubsection{\texorpdfstring{$\partial\log\beta/\partial\insrate$ and
$\partial\log\beta/\partial\delrate$}{partial log beta / partial lambda and partial log beta / partial mu}}

Using~\eqref{eq:dlogbeta-dl} with~\eqref{eq:singular-A}
and~\eqref{eq:singular-C}:
\begin{align*}
\frac{\partial\log\beta}{\partial\insrate}
  &= \frac{1}{\insrate}
     - \Bigl(\frac{1}{\varepsilon} + \frac{\evoltime}{2}\Bigr)
     + \Bigl(\frac{1}{\varepsilon} + \frac{s\evoltime}{2(1+s)}\Bigr)
     + O(\varepsilon)
\\
  &= \frac{1}{\delrate} - \frac{\evoltime}{2} + \frac{s\evoltime}{2(1+s)}
   = \frac{1}{\delrate} - \frac{\evoltime}{2(1+s)}
\end{align*}
Using~\eqref{eq:dlogbeta-dm} with~\eqref{eq:singular-B}
and~\eqref{eq:singular-D}:
\begin{align*}
\frac{\partial\log\beta}{\partial\delrate}
  &= \Bigl(\frac{1}{\varepsilon} - \frac{\evoltime}{2}\Bigr)
     - \Bigl(\frac{1}{\varepsilon} - \frac{s\evoltime}{2(1+s)}\Bigr)
     + O(\varepsilon)
\\
  &= -\frac{\evoltime}{2} + \frac{s\evoltime}{2(1+s)}
   = -\frac{\evoltime}{2(1+s)}
\end{align*}

Converting to log-elasticities:
\begin{equation}
\boxed{
\frac{\partial\log\beta}{\partial\log\insrate}
  \;\to\; \frac{s+2}{2(1+s)},
\qquad
\frac{\partial\log\beta}{\partial\log\delrate}
  \;\to\; -\frac{s}{2(1+s)}}
\label{eq:logbeta-elasticity-limit}
\end{equation}

\paragraph{Check.} As a consistency check, we verify that the sum of the
two elasticities (which gives the response to uniformly scaling both
rates at fixed $\kappa = \insrate/\delrate$) matches the direct
derivative of the limit formula.
The sum is $(s{+}2)/(2(1{+}s)) + (-s)/(2(1{+}s)) = 1/(1{+}s)$.
At the limit, $\beta = s/(1{+}s)$ where $s = \delrate\evoltime$,
so $d\log\beta/d\log s = 1/(1{+}s)$.  \checkmark

\subsubsection{\texorpdfstring{$\partial\log(1{-}\beta)/\partial\log\insrate$ and
$\partial\log(1{-}\beta)/\partial\log\delrate$}{partial log(1-beta) / partial log lambda and partial log(1-beta) / partial log mu}}

Using $\beta/(1-\beta) \to s$ at equal rates:
\begin{equation}
\boxed{
\frac{\partial\log(1{-}\beta)}{\partial\log\insrate}
  \;\to\; -\frac{s(s+2)}{2(1+s)},
\qquad
\frac{\partial\log(1{-}\beta)}{\partial\log\delrate}
  \;\to\; \frac{s^2}{2(1+s)}}
\label{eq:log1mbeta-elasticity-limit}
\end{equation}

\paragraph{Check.} Summing the two elasticities (uniform rate scaling):
$(-s^2{-}2s{+}s^2)/(2(1{+}s)) = -s/(1{+}s)$.
Direct: $d\log(1{-}\beta)/d\log s = -s/(1{+}s)$.  \checkmark

\subsubsection{\texorpdfstring{$\partial\log(1{-}\gamma)/\partial\log\insrate$ and
$\partial\log(1{-}\gamma)/\partial\log\delrate$}{partial log(1-gamma) / partial log lambda and partial log(1-gamma) / partial log mu}}

From~\eqref{eq:dlog1mgamma-dl}:
$\partial\log(1{-}\gamma)/\partial\insrate
= \partial\log\beta/\partial\insrate - 1/\insrate
\to 1/\delrate - \evoltime/(2(1{+}s)) - 1/\delrate
= -\evoltime/(2(1{+}s))$.

From~\eqref{eq:dlog1mgamma-dm}:
$\partial\log(1{-}\gamma)/\partial\delrate
= 1/\delrate + \partial\log\beta/\partial\delrate
  - \evoltime\alpha/\Phi
\to 1/\delrate - \evoltime/(2(1{+}s)) - \evoltime\alpha/\Phi$.

Converting:
\begin{equation}
\boxed{
\frac{\partial\log(1{-}\gamma)}{\partial\log\insrate}
  \;\to\; -\frac{s}{2(1+s)},
\qquad
\frac{\partial\log(1{-}\gamma)}{\partial\log\delrate}
  \;\to\; \frac{s+2}{2(1+s)} - \frac{s\alpha}{\Phi}}
\label{eq:log1mgamma-elasticity-limit}
\end{equation}

\paragraph{Check.} Summing the two elasticities (uniform rate scaling):
$-s/(2(1{+}s)) + (s{+}2)/(2(1{+}s)) - s\alpha/\Phi
= 1/(1{+}s) - s\alpha/\Phi$.
Direct: at the limit, $1-\gamma = s/((1{+}s)\Phi)$, so
$\log(1{-}\gamma) = \log s - \log(1{+}s) - \log\Phi$.
Under $s \to cs$:
$d/d\log c|_{c=1} = 1 - s/(1{+}s) - s\alpha/\Phi
= 1/(1{+}s) - s\alpha/\Phi$.  \checkmark

\subsubsection{\texorpdfstring{$\partial\log\gamma/\partial\log\insrate$ and
$\partial\log\gamma/\partial\log\delrate$}{partial log gamma / partial log lambda and partial log gamma / partial log mu}}

With $R = (1{-}\gamma)/\gamma$ at equal rates.
Since $1-\gamma = 1/((1{+}s)\phi)$:
\[
R = \frac{1}{(1+s)\phi - 1}
\]

\begin{equation}
\boxed{
\frac{\partial\log\gamma}{\partial\log\insrate}
  \;\to\; \frac{Rs}{2(1+s)},
\qquad
\frac{\partial\log\gamma}{\partial\log\delrate}
  \;\to\; -R\Bigl(\frac{s+2}{2(1+s)} - \frac{s\alpha}{\Phi}\Bigr)}
\label{eq:loggamma-elasticity-limit}
\end{equation}


\subsubsection{Summary Table: L'H\^opital Limits}
\label{sec:score-table-limits}

At $\insrate = \delrate$, with $s = \delrate\evoltime$,
$\alpha = e^{-s}$, $\Phi = 1-\alpha$,
$\phi = \Phi/s$, $R = 1/((1{+}s)\phi - 1)$:

\begin{center}
\renewcommand{\arraystretch}{1.8}
\begin{tabular}{c|cccccc}
& $\log\alpha$ & $\log(1{-}\alpha)$
& $\log\beta$ & $\log(1{-}\beta)$
& $\log\gamma$ & $\log(1{-}\gamma)$ \\
\hline
$\dfrac{\partial}{\partial\log\insrate}$
  & $0$
  & $0$
  & $\dfrac{s+2}{2(1+s)}$
  & $-\dfrac{s(s+2)}{2(1+s)}$
  & $\dfrac{Rs}{2(1+s)}$
  & $-\dfrac{s}{2(1+s)}$
\\[8pt]
$\dfrac{\partial}{\partial\log\delrate}$
  & $-s$
  & $\dfrac{s\alpha}{\Phi}$
  & $-\dfrac{s}{2(1+s)}$
  & $\dfrac{s^2}{2(1+s)}$
  & $-R\!\left(\dfrac{s{+}2}{2(1{+}s)}{-}\dfrac{s\alpha}{\Phi}\right)$
  & $\dfrac{s+2}{2(1+s)}{-}\dfrac{s\alpha}{\Phi}$
\\[8pt]
$\dfrac{\partial}{\partial\log(\insrate{+}\delrate)}$
  & $-s$
  & $\dfrac{s\alpha}{\Phi}$
  & $\dfrac{1}{1+s}$
  & $-\dfrac{s}{1+s}$
  & $R\!\left(\dfrac{s}{2(1{+}s)}{-}\dfrac{s{+}2}{2(1{+}s)}{+}\dfrac{s\alpha}{\Phi}\right)$
  & $\dfrac{1}{1+s}{-}\dfrac{s\alpha}{\Phi}$
\end{tabular}
\end{center}

\noindent
All entries are finite.
The third row equals the sum of the first two (as it must by
construction).
The $\log\alpha$ and $\log(1{-}\alpha)$ columns are trivially
continuous (no $\insrate$ dependence).
The $\log\beta$ and $\log(1{-}\beta)$ entries have a
particularly clean form involving only $s/(1{+}s)$.
The $\log\gamma$ entries additionally involve
$\alpha/\Phi = e^{-s}/(1-e^{-s}) = 1/(e^s - 1)$ and the
ratio $R = 1/((1{+}s)\phi - 1)$.


%%%%%%%%%%%%%%%%%%%%%%%%%%%%%%%%%%%%%%%%%%%%%%%%%%%%%%%%%%%%%%%%%%%%%%%%
\subsection{BDI Sufficient Statistics}
\label{sec:lhopital-bdi}

The EM-relevant quantities are the posterior expectations
$\mathbb{E}[B|i,j,T]$, $\mathbb{E}[D|i,j,T]$, $\mathbb{E}[S|i,j,T]$,
recovered from the score via the conservation law $i + B - D = j$:
\begin{align}
\mathbb{E}[S] &= \frac{j - i
  + \delrate\,\partial_\delrate\log P_{ij}
  - \insrate\,\partial_\insrate\log P_{ij}
  - \insrate\evoltime}{\insrate - \delrate}
\label{eq:ES-formula} \\
\mathbb{E}[B] &= \insrate\,\partial_\insrate\log P_{ij}
  + \insrate\,\mathbb{E}[S] + \insrate\evoltime
\label{eq:EB-formula} \\
\mathbb{E}[D] &= \delrate\,\partial_\delrate\log P_{ij}
  + \delrate\,\mathbb{E}[S]
\label{eq:ED-formula}
\end{align}

The score $\partial_\insrate\log P_{ij}$ itself is computed from
the transition count groups:
\[
\partial_\insrate\log P_{ij}
= \sum_{\xi \in \{\alpha,1{-}\alpha,\beta,\ldots\}}
  n_{\log\xi} \cdot \frac{\partial\log\xi}{\partial\insrate}
\]
Each derivative $\partial\log\xi/\partial\insrate$ has a
finite L'H\^opital limit (Section~\ref{sec:score-limits}),
so the score itself has a finite limit.
The denominator $\insrate - \delrate$ in~\eqref{eq:ES-formula}
is what creates the $0/0$ for $\mathbb{E}[S]$.

\paragraph{The $\mathbb{E}[S]$ singularity.}
The numerator of~\eqref{eq:ES-formula} also vanishes
at $\insrate = \delrate$, since the complete-data log-likelihood
for $\immrate = \insrate$ is
$\ell = B\log\insrate + D\log\delrate - (\insrate+\delrate)S
- \insrate\evoltime + c$,
and at $\insrate = \delrate$ the numerator becomes
$j - i + \mathbb{E}[D] - \mathbb{E}[B] = 0$ by conservation.

To resolve the $0/0$, write
$f(\insrate) = j - i + \delrate\,\partial_\delrate\log P
- \insrate\,\partial_\insrate\log P - \insrate\evoltime$
and $g(\insrate) = \insrate - \delrate$.
By L'H\^opital ($g' = 1$):
\begin{equation}
\lim_{\insrate\to\delrate} \mathbb{E}[S]
= f'(\delrate)
= \delrate\,\partial^2_{\insrate\delrate}\log P
  - \partial_\insrate\log P
  - \delrate\,\partial^2_{\insrate^2}\log P
  - \evoltime
\label{eq:ES-limit}
\end{equation}
where all derivatives are evaluated at $\insrate = \delrate$.

The first and second derivatives of $\log P$ that appear
in~\eqref{eq:ES-limit} can be expressed in terms of the
transition count groups and the L'H\^opital-limited score
derivatives from Section~\ref{sec:score-limits},
together with their first-order corrections in $\varepsilon$.
Specifically, the second derivatives of $\log P$ require
the $O(\varepsilon)$ terms in the Laurent expansions
\eqref{eq:singular-A}--\eqref{eq:singular-D}.

\paragraph{Alternative: direct gradient on $\delrate$.}
At equal rates, the BDI has a single parameter $\delrate$.
The gradient is:
\[
\frac{\partial\log P}{\partial\delrate}\bigg|_{\insrate=\immrate=\delrate}
= \sum_\xi n_{\log\xi}\,
  \Bigl(\frac{\partial\log\xi}{\partial\insrate}
  + \frac{\partial\log\xi}{\partial\delrate}\Bigr)
= \frac{\mathbb{E}[B{+}D]}{\delrate} - 2\mathbb{E}[S] - \evoltime
\]
This is well-defined using the L'H\^opital limits,
and gives the combined quantity
$\mathbb{E}[B{+}D] - \delrate(2\mathbb{E}[S]+\evoltime)$
without needing to separate $\mathbb{E}[B]$, $\mathbb{E}[D]$,
$\mathbb{E}[S]$ individually.
The M-step~\eqref{eq:mu-mle-equal} sets this to zero:
$\hat\delrate = (\mathbb{E}[B]+\mathbb{E}[D])/(2\mathbb{E}[S]+\evoltime)$.

To \emph{break out} of equal rates (determine whether
$\insrate$ should differ from $\delrate$), we need the
individual $\partial\log P/\partial\insrate$ and
$\partial\log P/\partial\delrate$, which the table in
Section~\ref{sec:score-table-limits} provides,
or we need to use additional information from the ancestral length distribution.


\subsection{Direct BDI Rate EM at Equal Rates}
\label{sec:equal-rate-em}

If one ignores the stationary root-length prior and considers only the conditional BDI bridge model with $\insrate = \immrate = \delrate$,
the complete-data log-likelihood~\eqref{eq:bdi-complete} reduces to:
\[
\ell = (B + D)\log\delrate - 2\delrate S - \delrate\evoltime + c
\]
This has a single free parameter $\delrate$.
Setting $\partial\ell/\partial\delrate = 0$:
\begin{equation}
\boxed{\hat{\delrate} = \frac{B + D}{2S + \evoltime}}
\label{eq:mu-mle-equal}
\end{equation}
where $B$, $D$, $S$ are understood as posterior expectations.

\paragraph{Implication.}
At equal rates, the M-step collapses to estimating a single rate
from the total event count and total sojourn time.


\subsection{\texorpdfstring{$\log\kappa$ and $\log(1{-}\kappa)$ Derivatives}{log kappa and log(1-kappa) Derivatives}}

The $\kappa = \insrate/\delrate$ and $(1{-}\kappa)$ groups
are not among the six $(\alpha,\beta,\gamma)$ factors but
appear in the transition matrix.
At $\kappa = 1$:
\begin{align*}
\frac{\partial\log\kappa}{\partial\insrate} &= \frac{1}{\insrate}
  \;\to\; \frac{1}{\delrate}, &
\frac{\partial\log\kappa}{\partial\delrate} &= -\frac{1}{\delrate}
\end{align*}
No singularity.  However, $\log(1{-}\kappa)$ diverges:
\[
\frac{\partial\log(1-\kappa)}{\partial\insrate}
= \frac{-\kappa}{(1-\kappa)\insrate}
\;\to\; -\infty
\]
Unlike the conditional WFST factors $\alpha,\beta,\gamma$, the stationary-length factors $\kappa$ and $1-\kappa$ do not admit a regular equal-rate limit in the full joint TKF model.
As $\kappa \to 1$, the equilibrium length distribution ceases to be normalizable, and the $\log(1-\kappa)$ contribution diverges.
Accordingly, the L'Hôpital limits derived above apply to the conditional bridge/WFST quantities, whereas the full stationary joint model must either remain in the subcritical regime $\kappa < 1$
or be reparameterized with a nonstationary root-length model.


% joint-vs-conditional.tex — Joint vs conditional likelihoods for TKF EM
% To be \input from lhopital-limits.tex

\subsection{Joint vs Conditional Likelihoods}
\label{sec:joint-vs-conditional}

The TKF91 transition matrix includes factors of $\kappa = \insrate/\delrate$
and $(1{-}\kappa)$ arising from the geometric$(\kappa)$ prior on
ancestor length: $\Pr(|\text{anc}| = i) = \kappa^i(1{-}\kappa)$.
This prior is part of the \emph{joint} likelihood $P(x,y)$
but can be factored out with respect to the ancestor-length prior to yield an \emph{ancestor length-conditioned} likelihood
$P_{\mathrm{cond}}(x,y \mid |x|=i)
= P(x,y) / [\kappa^i(1{-}\kappa)]$.

\paragraph{The four regimes.}
Of the four combinations of $\kappa < 1$ vs $\kappa = 1$
and joint vs ancestor length-conditioned, only one is degenerate:
\begin{center}
\begin{tabular}{lcc}
\toprule
& $\kappa < 1$ & $\kappa = 1$ \\
\midrule
Joint $P(x,y)$ & well-defined & \textbf{degenerate} ($\to 0$) \\
Ancestor length-conditioned $P_{\mathrm{cond}}$ & well-defined & well-defined \\
\bottomrule
\end{tabular}
\end{center}
At $\kappa = 1$, the geometric prior places vanishing mass on finite
sequences, so $P(x,y) \to 0$ and
$\log P \to -\infty$.
The ancestor length-conditioned WFST removes the $\kappa$ and $(1{-}\kappa)$
factors corresponding to the stationary length distribution from the transition matrix, leaving only
$(\alpha,\beta,\gamma)$-dependent terms.
These have finite L'H\^opital limits (Sections~\ref{sec:score-limits}--\ref{sec:lhopital-bdi}),
so $P_{\mathrm{cond}}$ is finite and well-defined even at $\kappa = 1$.

\paragraph{BDI sufficient statistics.}
The BDI sufficient statistics  $\mathbb{E}[B]$, $\mathbb{E}[D]$,
$\mathbb{E}[S]$ are most naturally recovered from the ancestor length-conditioned WFST, whose score depends only on the six TKF bridge groups 
$(\alpha, 1{-}\alpha, \beta, 1{-}\beta, \gamma, 1{-}\gamma)$. In the joint model, additional contributions from the ancestral length prior appear through the
$\kappa$ and $1{-}\kappa$ factors. These must be included in the M-step when optimizing the full stationary TKF likelihood, but they are conceptually distinct from the bridge statistics themselves.

Equations~(\ref{eq:exposure-tkf})-\ref{eq:deaths-tkf} give the correct expectations for the BDI sufficient statistics
in terms of the counts $\hat{n}_{ij}$ regardless of whether those counts were computed using the joint Pair HMM or the conditionally-normalized transducer,
as long as the WFST transition matrix $\tkfwfst$ is used for the counts $\emcount^\xi_{ij} = \xi \partial_\xi \log \tkfwfst_{ij}$
(e.g. using the tables in Section~\ref{sec:score-table-general} or Section~\ref{sec:score-table-limits}).
The M-step for the joint likelihood includes additional terms from the $\kappa$ and $1{-}\kappa$ factors, leading to a coupled system of equations that can be solved in closed form (Section~\ref{sec:bw-tkf91}).

\paragraph{Closed-form M-steps.}
Both formulations yield closed-form M-steps:
\begin{itemize}
\item \emph{Ancestor-length conditioned:} pure BDI exponential family.
  For the conditional model away from the $\kappa=1$ singularity of the joint prior,
  setting $\partial Q_{\mathrm{cond}}/\partial\insrate = 0$ gives
  \begin{align}
  \hat\insrate &= \frac{\max(0,\;\mathbb{E}[B] + \alpha_\insrate - 1)}
                       {\mathbb{E}[S] + \evoltime + \beta_\insrate}, &
  \hat\delrate &= \frac{\max(0,\;\mathbb{E}[D] + \alpha_\delrate - 1)}
                       {\mathbb{E}[S] + \beta_\delrate}
  \label{eq:mstep-cond}
  \end{align}
  where $(\alpha,\beta)$ are Gamma prior parameters.

\item \emph{Joint ($\kappa < 1$):}
  includes $L \log\kappa + M \log(1{-}\kappa)$.
  Since, for a single sequence-pair, $L  = i$ (ancestor length, exact) and
  $M  = 1$ (one End transition, exact),
  with corresponding accumulated values across multiple sequence-pairs,
  the optimality conditions
  $\partial Q/\partial\insrate = 0$ and $\partial Q/\partial\delrate = 0$
  form a coupled system in $(\insrate,\delrate)$.
  These can be reduced to a single quadratic equation in $\kappa$ with a closed-form solution, which can be substituted back to get $\hat\insrate$ and $\hat\delrate$,
  as described in Section~\ref{sec:bw-tkf91}.
\end{itemize}
In both cases, the M-step maximizes the correct
$Q$-function, so EM monotonicity is guaranteed for
the corresponding tracked objective
($\log P_{\mathrm{cond}}$ for the ancestor length-conditioned model, $\log P$ for the full joint model).


% irreversible.tex — Irreversible substitution models and κ > 1 regime
% To be \input from lhopital-limits.tex

%%%%%%%%%%%%%%%%%%%%%%%%%%%%%%%%%%%%%%%%%%%%%%%%%%%%%%%%%%%%%%%%%%%%%%%%
\subsection{\texorpdfstring{Irreversible Models and the $\insrate > \delrate$ Regime}{Irreversible Models and the lambda > mu Regime}}
\label{sec:irreversible}

\paragraph{The $\kappa > 1$ regime.}
When $\insrate > \delrate$, the ratio $\kappa = \insrate/\delrate > 1$
and the geometric$(\kappa)$ prior on ancestor length is not
normalizable: $\sum_{i=0}^\infty \kappa^i(1-\kappa)$ diverges.
Therefore the joint likelihood $P(x,y)$ does not exist.
However, the ancestor length-conditioned likelihood
$P_{\mathrm{cond}}(x,y \mid |x|)$
is well-defined for all $\kappa > 0$.
The transition parameters $\alpha$, $\beta$, $\gamma$ are
continuous functions of $(\insrate,\delrate,\evoltime)$
with no singularity at $\insrate = \delrate$ or $\insrate > \delrate$.
The BDI sufficient statistics $\mathbb{E}[B]$, $\mathbb{E}[D]$,
$\mathbb{E}[S]$ remain well-defined, and the conditioned
M-step~\eqref{eq:mstep-cond} applies unchanged:
it is unconstrained in $\insrate$ and $\delrate$ and permits
$\hat\insrate > \hat\delrate$ if the data support it.

\paragraph{Irreversible substitution models.}
For completeness, our implementations support dropping the
reversibility assumption for the substitution CTMC.
The integrals described in Section~\ref{sec:bridge-expectations} and in~\cite{HolmesRubin2002}
exploit the symmetric factorization of $\revsub$
using $\eqm$-weighted similarity, which yields real eigenvalues.
In the irreversible case, $\revsub$ is a general rate matrix
with no detailed-balance constraint.
The eigendecomposition $\revsub = V\,\operatorname{diag}(d)\,V^{-1}$
may have complex eigenvalues.
Assuming $\revsub$ is diagonalizable, the spectral formulae for the integrals of Section~\ref{sec:bridge-expectations} carry over,
with $V^{-1}$ replacing $V^\top$ and complex arithmetic throughout:
\[
I^{ab}_{ij}(\evoltime) =
\sum_{k,l} V_{ak}\,V^{-1}_{ki}\,
\frac{e^{d_k\evoltime} - e^{d_l\evoltime}}{d_k - d_l}\,
V_{jl}\,V^{-1}_{lb}
\]
with the convention $(e^{d_k\evoltime} - e^{d_l\evoltime})/(d_k - d_l) = \evoltime\,e^{d_k\evoltime}$
when $d_k = d_l$.
The result is real (complex eigenvalues come in conjugate pairs).

\paragraph{Fixed-point iteration for $\eqm$.}
In the reversible case, the equilibrium distribution $\eqm$ is
constrained to be the stationary distribution of $\revsub$ by
by detailed balance, and may either be fixed a priori or updated jointly with the reversible rate parameters during the M-step.
In the irreversible case, $\eqm$ must be recomputed as the
left eigenvector of $\revsub$ (with eigenvalue~0) after each
M-step update of $\revsub$ from the bridge-expectation sufficient statistics.
This gives an ECM (Expectation-Conditional-Maximization)
scheme~\cite{MengRubin1993}:
the E-step uses the current $(\revsub, \eqm)$;
the M-step updates $\revsub$, then $\eqm$ is recomputed from
the updated $\revsub$.
Standard monotonicity results apply providing each conditional maximization step is well-defined.

\paragraph{Differentiability.}

Standard autograd libraries provide general eigendecomposition and matrix-exponential primitives (with Fréchet derivatives).
The eigendecomposition route is convenient for EM-style closed-form updates, while gradient-based optimization is often more naturally based on the matrix exponential. We do not rely on differentiating through eigenvectors here.



\subsection{Summary of Limits}

\begin{center}
\begin{tabular}{lll}
\hline
\textbf{Quantity} & \textbf{Formula} & \textbf{Limit as $\insrate \to \delrate$} \\
\hline
$\alpha$ & $e^{-\delrate\evoltime}$ & $e^{-s}$ (continuous) \\
$\beta$ & $\frac{\insrate(\eta-\alpha)}{\delrate\eta - \insrate\alpha}$
  & $\frac{s}{1+s}$ \\
$1-\beta$ & & $\frac{1}{1+s}$ \\
$\gamma$ & $1 - \frac{\delrate\beta}{\insrate(1-\alpha)}$
  & $1 - \frac{1}{(1+s)\phi}$ \\
$\kappa$ & $\insrate/\delrate$ & $1$ \\
$\partial_{\log\insrate}\log\beta$
  & \eqref{eq:dlogbeta-dl}
  & $\frac{s+2}{2(1+s)}$ \\
$\partial_{\log\delrate}\log\beta$
  & \eqref{eq:dlogbeta-dm}
  & $-\frac{s}{2(1+s)}$ \\
$\hat{\delrate}_{\text{MLE}}$ & &
  $\frac{\mathbb{E}[B]+\mathbb{E}[D]}{2\mathbb{E}[S]+\evoltime}$ \\
\hline
\end{tabular}
\end{center}

where $s = \delrate\evoltime$ throughout.
